\section{Run the Page}
\label{sec:ch13-run-the-page}
\index{statement count!search page}%

The new root can express the request that \code{sessions} could not:

\begin{minted}{graphql}
query SearchPage($minimumScore: Float, $first: Int) {
  searchSessions(
    minimumScore: $minimumScore
    order: RATING_DESCENDING
    first: $first
  ) {
    nodes {
      session { id title }
      averageScore
      ratingCount
    }
    pageInfo { hasNextPage endCursor }
  }
}
\end{minted}

With \code{minimumScore} set to 4.0 and \code{first} set to 2, the router
returns:

\begin{minted}[fontsize=\footnotesize]{json}
{
  "data": {
    "searchSessions": {
      "nodes": [
        {
          "session": {
            "id": "U2Vzc2lvbjoy",
            "title": "Reading a Query Plan"
          },
          "averageScore": 4.5,
          "ratingCount": 2
        },
        {
          "session": {
            "id": "U2Vzc2lvbjox",
            "title": "Schemas That Outlive Their Authors"
          },
          "averageScore": 4,
          "ratingCount": 3
        }
      ],
      "pageInfo": {
        "hasNextPage": false,
        "endCursor": "MQ=="
      }
    }
  }
}
\end{minted}

The first row is now the 4.5 session. More important, \code{hasNextPage} is
about the filtered ordering, not the schedule ordering that happened to be
available elsewhere.

The Search log proves where the cut occurred. Its SQL filters by average score,
orders that column descending, then orders by session id before applying the
limit. The id tie-breaker is inside the same statement as the limit. It is not
a repair step in the router.

\begin{center}
\begin{tabular}{@{}lr@{}}
\toprule
Service & SQL statements \\
\midrule
Search & 1 \\
Sessions & 1 \\
Ratings & 0 \\
Speakers & 0 \\
\bottomrule
\end{tabular}
\end{center}

The Sessions statement is one batched entity lookup for the two chosen keys.
Ratings is absent because the averages used for discovery came from the Search
projection. Speakers is absent because the request did not select a speaker.

The plan has a \code{Sequence} with two children. The first is a
\code{Single} fetch from Search, subgraph id 3. The second is a
\code{BatchEntity} fetch from Sessions, subgraph id 0, and it declares that
it waits for fetch 0. Reversing those two steps would recreate the original
problem. The plan places them in the only order that can answer the request.

An unfiltered request is a separate completeness check. In rating-descending
order the four averages are 4.5, 4.0, 3.0, and null. The unrated session is
still discoverable and reports a rating count of 0. A projection built only by
grouping existing rating rows would silently lose it.

\subsection{The Consistency Bill}
\label{sec:ch13-consistency-bill}
\index{search projection!freshness}%

The deterministic seed isolates execution mechanics. A real system needs an
update path from schedule and rating changes into the projection. That might
be event-driven, a change-data-capture pipeline, or a periodic rebuild. The
choice belongs to the product's freshness requirement; Federation does not
choose it.

The API should say what that requirement permits. If a newly submitted rating
may take a minute to affect discovery, document and monitor that lag. If a
booking workflow requires read-your-write behavior, this eventually updated
projection may be the wrong mechanism for that workflow. A search index is not
a free join. It moves the join to write or ingestion time and makes staleness a
domain concern.

I also would not call the cursor stable merely because it has a deterministic
tie-breaker. The tie-breaker orders one snapshot. If scores change between two
requests, members can move across a cursor boundary. Snapshot tokens or a
versioned index are different features, and this chapter has not implemented
either.

That is the full trade: the graph gains a page it can honestly answer, and the
Search domain gains responsibility for completeness, freshness, and ordering.
